\documentclass{jsarticle}
\usepackage{amsmath}
\usepackage{amssymb}
\begin{document}
\title{}
\author{}
\date{}
\maketitle
   
  質量欠損の発見は、富山県の人が関わっている。
 重水素から $$ \Delta M = (A - Z)M_1 + ZM_1 $$ 
 と原子核の質量が重水素になると原子の質量が減少する。
 この減少するエネルギーから、特殊相対性理論の$ \Delta E \cong mc^2 $
 が発見された。遷移元素は、この典型元素の性質をランタノイドとアクチノイド
 で光量子仮説として表している。
 天秤による重さと質量の違いからこの重さをその対象とする原子の
 収率から質量がわかる。この天秤に例えられる３次元多様体による、
 ブラックホールとホワイトホールの５次元多様体のザイフェルト多様体から
 Higgs粒子との水素原子の質量とオイラーの定数の関係から、
 大域的微分方程式により、物体の質量の仕組みがわかる。
 $$ M_H = 1, 0.58, M = 1.58 $$
 から水素原子とオイラーの定数より、ゼータ関数を量子群と統合すると
 ヒッグス場の方程式となる。
 これらから、物体の収率から相対的な重さに密度和を積分すると質量が
 なぜ宇宙と異次元に存在しているかがわかる。
 $$ {d \over df}F = {d \over df}\int\int_M {1 \over (x \log x)^2}dx_m
  + {d \over df}\int\int_M {1 \over (y \log y)^{1 \over 2}}dy_m $$
 $$ {V(x) \over f(x)} = m(x) $$
 $$ E^2 = m^2c^4, R^{\mu\nu} - {1 \over 2}\Lambda g_{ij}
 = 8 \pi G T^{\mu\nu} $$
 $$ = 4 \pi G \rho $$
 $$ \Psi H = i \hbar \psi $$
 $$ {\Psi H \over i} = \hbar \psi $$
 $$ {{d \over df} \int \int_M { 1 \over (y \log y)^{1 \over 2}}dy_m
   \over {d \over df} \int \int_M { 1 \over (x \log x)^2}dx_m }
 = {1 \over i} $$


  $$ ||ds^2|| = e^{-2\pi T|\psi|}[\eta_{\mu\nu} + \bar{h}(x)]
   + T^2 d^2 \psi $$
   $$ T^2 d^2 \psi = [\pi], \sum_{k=0}^{\infty}a_k f^k(x) $$
   $$ C = \int\int_M {1 \over (x \log x)^2}dx_m $$
   $$ C = \int{1 \over n^s}dx - \log x $$
   $$ \lim_{x\to a}f(x) + \lim_{y\to b}f(y) = l + m $$
   $$ \log (x \log x) - 2(y \log y)^{1 \over 2} \ge 0 $$
   $$ {d \over dx}\int^{x}_a f(x)dx = f(a) $$
   $$ M_H = 1 + C(オイラー定数), 
    \int f(z)dz = z + 1 $$
    $$ \mathcal{O}(x) = e^{-2 \pi T|\psi|} $$



 


\end{document}
